\SourceLeaf{250}%
und \(q\) den Winkel, den das erste Element dieser Linie mit dem ersten Element einer gegebenen, von \(A\) ausgehenden kürzesten Linie bildet. Es sei \(B\) ein fester Punkt auf dieser Linie, auf der \(q=0\) gilt, und \(C\) ein anderer fester Punkt der Fläche, für den wir den Wert von \(q\) einfach mit \(A\) bezeichnen. Nehmen wir an, die Punkte \(B,C\) seien durch eine kürzeste Linie verbunden. Die von \(B\) aus gemessenen Strecken dieser Linie bezeichnen wir wie in Artikel 18 allgemein mit \(s\). Ebenso bezeichnen wir mit \(\theta\) den Winkel, den ein beliebiges Element \(\mathrm{d}s\) mit dem Element \(\mathrm{d}p\) bildet. Schließlich seien \(\theta^0,\theta'\) die Werte des Winkels \(\theta\) in den Punkten \(B,C\). Wir haben also auf der gekrümmten Fläche ein von kürzesten Linien begrenztes Dreieck. Seine Winkel bei \(B\) und \(C\) bezeichnen wir einfach mit denselben Buchstaben. Der erste ist die Ergänzung des Winkels \(\theta^0\) zu \(180^0\), der zweite der Winkel \(\theta'\) selbst. Wer unsere Rechnung betrachtet, sieht leicht, dass alle Winkel in Zahlen anstelle von Graden ausgedrückt sind und der Winkel \(57^0\,17'\,45''\), zu dem ein dem Radius gleicher Bogen gehört, als Einheit gilt. Bezeichnen wir also den Kreisumfang mit \(2\pi\), so müssen wir setzen
\[
\theta^0=\pi-B,\qquad\theta'=C
\]
Untersuchen wir nun die Gesamtkrümmung dieses Dreiecks. Sie ist \(=\int k\mathrm{d}\sigma\), wobei \(\mathrm{d}\sigma\) ein Flächenelement des Dreiecks bezeichnet. Da dieses Element durch \(m\mathrm{d}p\mathord{.}\mathrm{d}q\) ausgedrückt wird, müssen wir das Integral \(\int\!\int km\mathrm{d}p\mathord{.}\mathrm{d}q\) über die ganze Dreiecksfläche bestimmen. Beginnen wir mit der Integration nach \(p\). Wegen \(k=-\frac{1}{m}\mathord{.}\frac{\mathrm{d}\mathrm{d}m}{\mathrm{d}p^2}\) erhalten wir \(\mathrm{d}q\mathord{.}(\mathrm{Const.}-\frac{\mathrm{d}m}{\mathrm{d}p})\) für die Gesamtkrümmung des Flächenstücks zwischen den Linien des ersten Systems, denen die Werte \(q,q+\mathrm{d}q\) der zweiten Veränderlichen entsprechen. Da diese Krümmung für \mbox{\(p=0\)} verschwinden muss, muss die durch die Integration eingeführte Konstante gleich dem Wert von \(\frac{\mathrm{d}m}{\mathrm{d}p}\) für \(p=0\) sein, also gleich eins. Wir erhalten somit \(\mathrm{d}q(1-\frac{\mathrm{d}m}{\mathrm{d}p})\), wobei für \(\frac{\mathrm{d}m}{\mathrm{d}p}\) der Wert am Ende dieses Flächenstücks auf der Linie \(CB\) einzusetzen ist. Auf dieser Linie gilt jedoch nach dem vorhergehenden Artikel \(\frac{\mathrm{d}m}{\mathrm{d}p}\mathord{.}\mathrm{d}q=-\mathrm{d}\theta\), sodass unser Ausdruck in \(\mathrm{d}q+\mathrm{d}\theta\) übergeht. Führen wir nun die zweite Integration von \(q=0\) bis \(q=A\) aus, so erhalten wir die Gesamtkrümmung des Dreiecks \(=A+\theta'-\theta^0=A+B+C-\pi\).

Die Gesamtkrümmung ist gleich dem Flächeninhalt des Teils der Kugeloberfläche, der dem Dreieck entspricht, mit positivem oder negativem Vorzeichen, je nachdem, ob die gekrümmte Fläche, auf der das Dreieck liegt, konkav-konkav oder konkav-konvex ist. Als Flächeneinheit dient das Quadrat, dessen Seite die Längeneinheit (der Radius der Kugel) ist. Damit ist der Flächeninhalt der ganzen Kugeloberfläche \(=4\pi\). Der Teil der Kugeloberfläche,
